\documentclass[11pt,a4paper]{article}
\usepackage[margin=2.5cm]{geometry}
\usepackage{amsmath,amssymb}
\usepackage{booktabs}
\usepackage{parskip}
\usepackage{hyperref}
\usepackage{xcolor}
\usepackage{graphicx}
\usepackage{siunitx}

\hypersetup{colorlinks=true,linkcolor=blue,citecolor=blue,urlcolor=blue}

\newcommand{\DS}{\Delta\Sigma}
\newcommand{\Msun}{\,M_{\odot}}
\newcommand{\Mpc}{\,\mathrm{Mpc}}
\newcommand{\unitvec}[1]{\hat{\mathbf{#1}}}
\newcommand{\Reff}{r_{\rm eff}}

\title{\textbf{Simulating Triaxial NFW Clusters and Their Lensing Signal}\\[0.4em]
\large \texttt{xlensing.sims} -- algorithmic and mathematical description}
\author{destacking project}
\date{2026 June 9}

\begin{document}
\maketitle

\begin{abstract}
We document the algorithm used by \texttt{xlensing.sims} to inject the
weak-lensing signal of a triaxial NFW cluster into a background-galaxy
catalogue. The pipeline has three steps: (i) draw a 3D triaxial ellipsoid
with CDM-like axis-ratio priors and an isotropic line of sight, (ii)
project the ellipsoid to a 2D ellipse via the Schur complement of its
shape matrix, and (iii) replace the spherical NFW radius by an
effective elliptic radius that preserves the area of isodensity contours,
evaluating the standard NFW shear of Wright \& Brainerd (1999) at that
effective radius with the shear direction taken from the gradient of the
elliptic-radius field. The whole construction reduces analytically and
exactly to the spherical case in the limit $q_{\rm maj} = q_{\rm min}$.
\end{abstract}

\tableofcontents
\bigskip


\section{Overview and notation}

We use the following symbols throughout.

\begin{tabular}{ll}
$z_l$, $M_{200m}$, $c_{200m}$    & cluster lens redshift, mass, NFW concentration \\
$z_s$                            & background source redshift \\
$D_A(z_1, z_2)$                  & angular diameter distance between $z_1$ and $z_2$ \\
$\bar{\rho}_M(z)$                & mean matter density at $z$ \\
$\Sigma_{\rm crit}(z_l, z_s)$    & lensing critical surface mass density \\
$r_{200m}$, $r_s$                & virial and scale radii of the NFW profile \\
$\delta_c$                       & NFW characteristic over-density \\
$g_{\rm NFW}(x)$                 & dimensionless NFW shear profile (Wright \& Brainerd 1999) \\
$(a, b, c)$                      & 3D semi-axes of the triaxial ellipsoid, $a \ge b \ge c$ \\
$(\theta_{\rm los}, \phi_{\rm los})$ & line-of-sight direction in the ellipsoid frame \\
$(q_{\rm maj}, q_{\rm min}, \psi)$ & projected semi-major / semi-minor axes
                                     and position angle \\
$\Reff$                          & elliptic effective radius (Section~\ref{sec:effr}) \\
$\gamma_t$                       & tangential NFW shear at a given radius \\
$g = g_1 + i g_2$                & complex shear in the (RA, Dec) sky frame \\
\end{tabular}

\bigskip

The algorithm is the following loop, for each cluster:
\begin{enumerate}
\item Draw $(a, b, c, \theta_{\rm los}, \phi_{\rm los})$
      (Section~\ref{sec:draw}).
\item Project the ellipsoid to a sky-plane ellipse
      $(q_{\rm maj}, q_{\rm min}, \psi)$
      (Section~\ref{sec:proj}).
\item For every source in the catalogue, compute the elliptic
      effective radius $\Reff$ and shear orientation $\phi_{\rm eff}$,
      then apply the spherical NFW shear at $\Reff$ rotated by $\phi_{\rm eff}$
      (Sections~\ref{sec:effr}--\ref{sec:apply}).
\end{enumerate}

The code lives in \texttt{xlensing.sims}; the relevant functions are
\texttt{sample\_triaxial\_shapes}, \texttt{project\_ellipsoid},
\texttt{triaxial\_NFW\_wcs\_cluster\_shear}, and
\texttt{apply\_triaxial\_NFW\_shear\_region}.

\section{Spherical NFW shear (reference)}

The triaxial pipeline reuses the spherical NFW shear of
Wright \& Brainerd (1999) at an elliptic effective radius. The spherical
expressions used in the simulator are:
\begin{align}
\bar\rho_M(z) &= \Omega_M\,\rho_{\rm crit}(z), \\[2pt]
r_{200m}(z_l, M) &= \left[\frac{3 M}{4\pi \Delta_{m}\bar\rho_M(z_l)}\right]^{1/3},
   \quad \Delta_m = 200, \\[2pt]
\delta_c(c) &= \frac{\Delta_m}{3} \cdot \frac{c^3}{\ln(1+c) - c/(1+c)}, \\[2pt]
\Sigma_{\rm crit}(z_l, z_s) &= \frac{c_{\rm light}^2}{4\pi G}\,
   \frac{D_A(0, z_s)}{D_A(0, z_l)\,D_A(z_l, z_s)}, \\[2pt]
r_s(z_l, M, c) &= r_{200m}(z_l, M)/c.
\end{align}
The dimensionless NFW shear $g_{\rm NFW}(x)$ with $x = r/r_s$ is the
closed-form expression of Wright \& Brainerd (their eq.~14):
\begin{equation}
g_{\rm NFW}(x) =
\begin{cases}
\displaystyle
   \frac{8\,\mathrm{arctanh}\!\sqrt{(1-x)/(1+x)}}{x^2\sqrt{1-x^2}}
 + \frac{4}{x^2}\,\ln(x/2)
 - \frac{2}{x^2-1}
 + \frac{4\,\mathrm{arctanh}\!\sqrt{(1-x)/(1+x)}}{(x^2-1)\sqrt{1-x^2}},
   & x < 1,\\[10pt]
\displaystyle
   \frac{8\,\mathrm{arctan}\!\sqrt{(x-1)/(1+x)}}{x^2\sqrt{x^2-1}}
 + \frac{4}{x^2}\,\ln(x/2)
 - \frac{2}{x^2-1}
 + \frac{4\,\mathrm{arctan}\!\sqrt{(x-1)/(1+x)}}{(x^2-1)\sqrt{x^2-1}},
   & x > 1,\\[10pt]
\displaystyle \frac{10}{3} + 4\ln(1/2),
   & x = 1.
\end{cases}
\end{equation}
The spherical tangential shear at a projected radius $r$ is then
\begin{equation}
\gamma_t^{\rm sph}(r;\,z_l, M, c, z_s)
 \;=\; \frac{r_s\,\delta_c\,\bar\rho_M(z_l)}{\Sigma_{\rm crit}(z_l, z_s)}
       \;g_{\rm NFW}\!\left(\frac{r}{r_s}\right).
\label{eq:gammat-sph}
\end{equation}
For a spherical halo, the complex shear in the equatorial frame is
$g = \gamma_t^{\rm sph}(r)\,e^{2 i \theta}$ where $\theta$ is the polar angle
of the source measured from north toward east.

\section{Drawing the 3D ellipsoid}
\label{sec:draw}

\paragraph{Axis ratios.}
For each cluster $i$ we draw two axis ratios from priors chosen to mimic
CDM N-body shape distributions (Allgood et al.~2006; Schneider, Frenk
\& Cole 2012):
\begin{align}
b_i/a_i &\sim \mathcal{U}(q_{\rm ba,\,lo},\;q_{\rm ba,\,hi}),
   \quad \text{default } (q_{\rm ba,\,lo}, q_{\rm ba,\,hi}) = (0.5, 1.0),\\
c_i/a_i &\sim \mathcal{U}(q_{\rm ca,\,lo},\;\,b_i/a_i),
   \quad \text{default } q_{\rm ca,\,lo} = 0.3.
\end{align}
The hierarchical bound $c/a \le b/a$ enforces the ordering
$a \ge b \ge c$ exactly. The semi-major axis $a_i$ is set to the spherical
virial radius $r_{200m}(z_i, M_i)$; the other two semi-axes are
$b_i = (b/a)_i\, a_i$ and $c_i = (c/a)_i\, a_i$.

\paragraph{Line of sight.}
The orientation of the principal axes with respect to the observer is
drawn isotropically. Writing the LOS unit vector
in the ellipsoid principal frame as
\begin{equation}
\unitvec{n}(\theta_{\rm los}, \phi_{\rm los})
 \;=\; (\sin\theta_{\rm los}\cos\phi_{\rm los},\;
        \sin\theta_{\rm los}\sin\phi_{\rm los},\;
        \cos\theta_{\rm los}),
\end{equation}
the prior is uniform on the sphere:
\begin{equation}
\cos\theta_{\rm los} \sim \mathcal{U}(0, 1),\qquad
\phi_{\rm los}       \sim \mathcal{U}(0, 2\pi).
\end{equation}
Drawing the upper half-sphere ($\cos\theta_{\rm los} > 0$) loses no
generality because the ellipsoid is centrally symmetric. The
\verb|isotropic_los=False| option freezes $\theta_{\rm los}=0$ to
produce face-on views along the $c$-axis (useful for diagnostics).

\section{Projecting the ellipsoid onto the sky}
\label{sec:proj}

The 3D ellipsoid surface $\mathbf{x}^\top Q\, \mathbf{x} = 1$ has the
diagonal shape matrix
\begin{equation}
Q \;=\; \mathrm{diag}\!\left(\tfrac{1}{a^2},\;\tfrac{1}{b^2},\;\tfrac{1}{c^2}\right)
\end{equation}
in the principal frame. To project it along $\unitvec{n}$ onto the
plane perpendicular to $\unitvec{n}$ (the sky plane), we eliminate the
LOS component using the Schur complement of $Q$:
\begin{equation}
M_{ij} \;=\; \unitvec{e}_i^\top Q\,\unitvec{e}_j
            \;-\; \frac{(\unitvec{e}_i^\top Q\,\unitvec{n})\,
                          (\unitvec{e}_j^\top Q\,\unitvec{n})}
                        {\unitvec{n}^\top Q\,\unitvec{n}},
\qquad i,j \in \{1, 2\},
\label{eq:Mij}
\end{equation}
where $(\unitvec{e}_1, \unitvec{e}_2)$ is any orthonormal basis of the
sky plane. We build it by Gram--Schmidt:
\begin{equation}
\unitvec{e}_1 \;=\;
   \frac{\unitvec{a} - (\unitvec{a}\cdot\unitvec{n})\,\unitvec{n}}
        {\|\unitvec{a} - (\unitvec{a}\cdot\unitvec{n})\,\unitvec{n}\|},
\qquad
\unitvec{e}_2 \;=\; \unitvec{n} \times \unitvec{e}_1,
\end{equation}
where $\unitvec{a}=(1,0,0)$ is the principal $a$-axis (with a fallback
to $(0,1,0)$ when the LOS is too close to the $a$-axis to keep this
construction numerically well-conditioned). The position angle of the
projected major axis is therefore measured from $\unitvec{e}_1$, i.e.\
from the sky-plane projection of the cluster's $a$-axis.

\paragraph{Derivation.}
A direction $\mathbf{x}$ in 3D can be decomposed as
$\mathbf{x} = u\,\unitvec{e}_1 + v\,\unitvec{e}_2 + w\,\unitvec{n}$.
The ellipsoid equation becomes
\begin{equation}
1 = u^2 Q_{11} + v^2 Q_{22} + w^2 (\unitvec{n}^\top Q \unitvec{n})
   + 2 u v Q_{12} + 2 u w (\unitvec{e}_1^\top Q \unitvec{n})
   + 2 v w (\unitvec{e}_2^\top Q \unitvec{n}),
\end{equation}
with $Q_{ij} \equiv \unitvec{e}_i^\top Q\, \unitvec{e}_j$ for the
sky-plane block. Fixing $(u, v)$ and asking for the projection (the
boundary of the shadow on the sky plane), $w$ takes the value that
minimises the LHS over the LOS. The minimum is reached at
$w^* = -[u(\unitvec{e}_1^\top Q\unitvec{n}) + v(\unitvec{e}_2^\top Q\unitvec{n})]/(\unitvec{n}^\top Q\unitvec{n})$,
and substituting back yields $\mathbf{u}^\top M\, \mathbf{u} = 1$ with
$M$ exactly as in \eqref{eq:Mij}. This is the standard projection-by-
Schur-complement result for an ellipsoid.

\paragraph{Diagonalising the projection.}
Diagonalising the $2\times2$ symmetric $M$ gives the projected
semi-axes and orientation:
\begin{equation}
M\,\mathbf{v}_k = \lambda_k\,\mathbf{v}_k, \quad k = 1, 2 \text{ (with }\lambda_1 \le \lambda_2\text{)},
\qquad
q_{\rm maj} = \frac{1}{\sqrt{\lambda_1}},\quad
q_{\rm min} = \frac{1}{\sqrt{\lambda_2}},
\end{equation}
\begin{equation}
\psi \;=\; \arctan2(\mathbf{v}_1[2],\;\mathbf{v}_1[1]),
\end{equation}
with $\psi$ measured counter-clockwise from $\unitvec{e}_1$. A fixed
sky offset can be added a posteriori to align the principal frame with
the (RA, Dec) frame.

\section{Effective-radius mapping}
\label{sec:effr}

Given a projected ellipse $(q_{\rm maj}, q_{\rm min}, \psi)$ on the sky
plane, every source at a 2D offset $(\Delta x, \Delta y)$ (east, north)
from the cluster centre is mapped to an elliptic effective radius. The
construction is the standard ``circularised'' radius that preserves the
area of isodensity contours:
\begin{align}
u &=  \Delta x \sin\psi + \Delta y \cos\psi,
   \quad \text{(along the projected major axis)}\\
v &=  \Delta x \cos\psi - \Delta y \sin\psi,
   \quad \text{(along the projected minor axis)}\\
\Reff &=  \sqrt{\frac{u^2}{q_{\rm maj}^2} + \frac{v^2}{q_{\rm min}^2}}
         \;\sqrt{q_{\rm maj}\,q_{\rm min}}.
\label{eq:reff}
\end{align}
The first factor is the dimensionless elliptic radius
$\xi=\sqrt{(u/q_{\rm maj})^2 + (v/q_{\rm min})^2}$; the second factor
$\sqrt{q_{\rm maj}\,q_{\rm min}}$ converts $\xi$ back to a physical
length while keeping the equal-area normalisation. In the limit
$q_{\rm maj}=q_{\rm min}\equiv q$ this collapses to
$\Reff = q\,(\sqrt{u^2+v^2})/q = \sqrt{u^2+v^2}$, the spherical projected
radius, exactly. This guarantees the simulator is consistent with the
spherical NFW case as the ellipticity is taken to zero.

\paragraph{Shear orientation.}
A spherical NFW shear at radius $r$ is purely tangential to the
projected radius vector. For an elliptic surface density, the shear at
$(\Delta x, \Delta y)$ is tangent to the local isodensity contour,
i.e.\ perpendicular to the gradient $\nabla \Reff$ of the
effective-radius field. The components of $\nabla \Reff$, rotated back to
the (east, north) frame, are:
\begin{align}
g_x &= \sin\psi\,\frac{u}{q_{\rm maj}^2}
     + \cos\psi\,\frac{v}{q_{\rm min}^2},
\\
g_y &= \cos\psi\,\frac{u}{q_{\rm maj}^2}
     - \sin\psi\,\frac{v}{q_{\rm min}^2},
\\
\phi_{\rm eff} &= \arctan2(g_x,\; g_y).
\label{eq:phieff}
\end{align}
The angle $\phi_{\rm eff}$ is the position angle (from north toward east)
of the local isodensity normal, and the shear is rotated by this angle
on the sky. Again, $q_{\rm maj}=q_{\rm min}$ gives
$\phi_{\rm eff}=\arctan2(\Delta x, \Delta y)=\theta$, recovering the
spherical convention.

\section{The triaxial NFW shear}

The complex shear at a source position $(\Delta x, \Delta y)$ relative
to the cluster centre, for a cluster of mass $M$, concentration $c$,
redshift $z_l$, and source redshift $z_s$, is
\begin{equation}
\boxed{\,
\gamma(\Delta x, \Delta y;\,z_l, M, c, z_s,\, q_{\rm maj}, q_{\rm min}, \psi)
\;=\;
\gamma_t^{\rm sph}\!\left(\Reff;\,z_l, M, c, z_s\right)\;
e^{2i\,\phi_{\rm eff}},
\,}
\label{eq:tri-shear}
\end{equation}
with $\gamma_t^{\rm sph}$ from \eqref{eq:gammat-sph} evaluated at the
effective radius \eqref{eq:reff} and the orientation $\phi_{\rm eff}$
from \eqref{eq:phieff}. The implementation uses
$\max(\Reff, 10^{-4}\,r_s)$ as the argument of $g_{\rm NFW}$ to avoid
the singularity at $x \to 0$.

\paragraph{Approximations baked into the model.}
Equation~\eqref{eq:tri-shear} is an \emph{affine-rescaling} model: it
treats the triaxial halo as a spherical one whose isodensity contours
are stretched into ellipses on the sky. The radial profile is left
unchanged. This neglects (i) the radial variation of the projected
ellipticity that a true 3D NFW would have (the inner ellipticity is
slightly larger than at large radii); (ii) the fact that the projected
surface density of a true triaxial NFW is not strictly an ellipse with
the same axis ratio as the dark-matter ellipsoid; (iii) any radial
dependence of the principal-axis orientation. For weak-lensing
applications around $R \sim r_s$ to a few times $r_{200m}$, these are
sub-leading: the dominant effect is the LOS-driven amplitude bias
captured by the projected ellipse axis ratio, which is what
\eqref{eq:tri-shear} retains by construction.

\section{Applying the shear to a galaxy catalogue}
\label{sec:apply}

The function \texttt{apply\_triaxial\_NFW\_shear\_region} takes a
cluster description and a galaxy catalogue and rewrites the galaxies'
ellipticities to include the cluster's shear. Each cluster row carries
$[\,\mathrm{RA}, \mathrm{Dec}, z, M_{200}, c_{200}, a, b, c,
 \theta_{\rm los}, \phi_{\rm los},\, (\psi_{\rm offset})\,]$. The
projection $(q_{\rm maj}, q_{\rm min}, \psi_{\rm intr})$ is computed
once per cluster from $(a, b, c, \theta_{\rm los}, \phi_{\rm los})$ via
\texttt{project\_ellipsoid}; the final PA is
$\psi = \psi_{\rm intr} + \psi_{\rm offset}$, where the optional offset
rotates the projected major axis into the (RA, Dec) frame.

For every galaxy $j$:
\begin{enumerate}
\item Apply the foreground cut: only sources with $z_j > z_l$ are
      sheared (the rest carry $\gamma_j = 0$).
\item Compute the angular separation $\Delta\theta_j$ and azimuth
      $\vartheta_j$ from the cluster centre with Vincenty's formula
      (\texttt{equatorial\_to\_polar}).
\item Convert to the physical offset at the lens plane:
\begin{equation}
\Delta x_j = D_A(0, z_l)\,\Delta\theta_j\sin\vartheta_j,\qquad
\Delta y_j = D_A(0, z_l)\,\Delta\theta_j\cos\vartheta_j.
\end{equation}
\item Compute $R_{{\rm eff},j}$ and $\phi_{{\rm eff},j}$ from
      \eqref{eq:reff}--\eqref{eq:phieff}.
\item Compute the critical surface density
      $\Sigma_{\rm crit}(z_l, z_j)$ and the spherical tangential shear
      $\gamma_t^{\rm sph}(R_{{\rm eff},j})$ from \eqref{eq:gammat-sph} using
      $r_s = r_{200m}(z_l, M)/c$ and the cluster's $\delta_c$.
\item Form $\gamma_j = \gamma_t^{\rm sph}(R_{{\rm eff},j})\,e^{2 i \phi_{{\rm eff},j}}$.
\item Add the shear to the source ellipticity using the standard
      ellipticity-addition rule
      $e^{\rm out}_j = (e^{\rm in}_j + \gamma_j)/(1 + \gamma_j^* e^{\rm in}_j)$,
      implemented in \texttt{add\_shears}.
\end{enumerate}

\section{Drawing a population of triaxial clusters}

The helper \texttt{sample\_triaxial\_shapes(z, M, ba\_range, ca\_lo,
isotropic\_los, rng)} returns, for each input $(z_i, M_i)$, the
quantities used downstream:
$(q_{\rm ba}, q_{\rm ca}, \theta_{\rm los}, \phi_{\rm los},
r_{\rm vir}, q_{\rm maj}, q_{\rm min}, \psi_{\rm intr})$.
A typical injection-recovery experiment then calls
\texttt{apply\_triaxial\_NFW\_shear\_region} for each cluster, accumulating
the perturbation on the same shear catalogue. The mean residual at fixed
$(M, c, z)$ measures the amplitude bias of the spherical NFW
template when applied to triaxial systems projected at random
orientations: at $z\!\sim\!0.4$, $M\!\sim\!10^{14}\Msun$, $c\!\sim\!5$
and the default priors, the median spherical-vs-triaxial $\DS$ deficit
is $\sim$15--20\,\% and is approximately constant in $R$ -- this is the
``triaxiality bias'' that motivates the amplitude-correction parameter
$A$ in the destacking pipeline.

\section{Sanity checks and limits}

\begin{itemize}
\item \textbf{Spherical limit.} For $q_{\rm maj} = q_{\rm min}$,
      $\Reff = \sqrt{\Delta x^2 + \Delta y^2}$ and
      $\phi_{\rm eff} = \arctan2(\Delta x, \Delta y) = \vartheta$.
      Equation~\eqref{eq:tri-shear} then reproduces the spherical-NFW
      shear of \texttt{NFW\_wcs\_cluster\_shear} exactly.
\item \textbf{Area preservation.} The ellipse
      $u^2/q_{\rm maj}^2 + v^2/q_{\rm min}^2 = 1$ has area
      $\pi q_{\rm maj}q_{\rm min}$. The choice
      $\Reff = \xi\sqrt{q_{\rm maj}q_{\rm min}}$ ensures that the
      $\Reff = R_0$ contour bounds the same area
      $\pi R_0^2$ as the spherical contour at radius $R_0$.
\item \textbf{Diagonalisation continuity.} The eigendecomposition of
      $M$ has a sign ambiguity in the eigenvectors; we resolve it by
      sorting eigenvalues ascending and reading the major-axis
      orientation off $\mathbf{v}_1$. The position angle returned by
      \texttt{project\_ellipsoid} can therefore differ from a previous
      implementation's by $\pi$ -- harmless because
      $\gamma e^{2 i (\phi+\pi)} = \gamma e^{2 i \phi}$.
\item \textbf{Numerical guards.} The Gram--Schmidt direction $\unitvec{e}_1$
      becomes ill-conditioned when the LOS is within
      $\sim 10^{-10}\,{\rm rad}$ of the $a$-axis. In that case the code
      uses $\unitvec{a}=(0,1,0)$ as a fallback. The NFW shear is
      evaluated at $\max(\Reff, 10^{-4}\,r_s)$ to avoid the inner-radius
      singularity.
\item \textbf{LOS-only triaxiality bias.} Because the radial profile is
      not modified, the bias produced by this simulator on a stacked
      $\DS$ relative to a spherical NFW is dominated by the
      $\sqrt{q_{\rm maj}\,q_{\rm min}}$ rescaling of the radius and the
      $\Sigma_{\rm crit}$ amplitude. The simulator is therefore an
      \emph{LOS-projection-only} model: it captures the dominant
      systematic but is not a substitute for stacking a 3D triaxial NFW
      with a proper triaxial $\Sigma$ template at $r \ll r_s$.
\end{itemize}

\section{Code anchors}

\begin{tabular}{ll}
\verb|sample_triaxial_shapes|         & Sections~\ref{sec:draw}--\ref{sec:proj} \\
\verb|project_ellipsoid|              & Section~\ref{sec:proj} \\
\verb|triaxial_NFW_wcs_cluster_shear| & Sections~\ref{sec:effr} and \eqref{eq:tri-shear} \\
\verb|apply_triaxial_NFW_shear_region| & Section~\ref{sec:apply} \\
\verb|NFW_tangential_shear|, \verb|gNFW| & Equation~\eqref{eq:gammat-sph} \\
\verb|equatorial_to_polar|            & Vincenty separation/azimuth, step 2 of §\ref{sec:apply} \\
\verb|add_shears|                     & ellipticity addition, step 7 of §\ref{sec:apply} \\
\end{tabular}

\end{document}
